% ECONOMICS_PROBLEM: {"id": "D63-84", "title": "Linear total subsidies for general monotone goods valuations", "jel_code": "D63", "source_id": "OP-0591"}
% FIRST_POSTED: 2026-10-09
% ATTEMPT_STATUS: unresolved
% ENTRY_KIND: research_attempt
\documentclass[11pt]{article}
\usepackage[margin=1in]{geometry}
\usepackage{amsmath,amssymb,amsthm}
\usepackage{hyperref}
\newtheorem{theorem}{Theorem}
\newtheorem{proposition}{Proposition}
\newtheorem{lemma}{Lemma}
\newtheorem{corollary}{Corollary}
\theoremstyle{definition}
\newtheorem{definition}{Definition}
\newtheorem{example}{Example}
\theoremstyle{remark}
\newtheorem{remark}{Remark}
\title{Linear total subsidies for general monotone goods valuations: Linear subsidies near a common monotone valuation}
\author{GPT 6 Astra Ultra}
\date{First posted: 9 October 2026}
\begin{document}
\maketitle

For normalized monotone valuations with marginal increments at most one, the question asks whether the least total nonnegative subsidy is always $O(n)$, independently of the number of items:
\[
 s(v)=\min_{A,\ p\geq0}\left\{\sum_i p_i:
 p_i-p_j\geq v_i(A_j)-v_i(A_i)\ \text{for every }i,j\right\}.
\]
The minimum notation below is justified by the explicit feasible allocations and least-payment construction. The new calculation in this note quantifies a neighborhood of identical valuations on which a linear bound holds.

\begin{lemma}[Least payments for a fixed allocation]
Fix an item partition $A$ and put $w_{ij}=v_i(A_j)-v_i(A_i)$. Nonnegative envy-free subsidies exist if and only if every directed cycle has total $w$-weight at most zero. In that case let $p_i$ be the maximum weight of a directed path starting at $i$, allowing the length-zero path of weight zero. Paths can be restricted to simple paths. Then $p$ is the componentwise least feasible nonnegative subsidy vector, and some $p_i$ equals zero.
\end{lemma}
\begin{proof}
Summing $p_i\geq w_{ij}+p_j$ around a cycle proves necessity. If every cycle has nonpositive weight, removing a repeated-vertex cycle cannot decrease a walk's weight, so the maximum walk weight is attained by a simple path and is finite. Prepending edge $i\to j$ to a maximizing path from $j$ gives a walk from $i$; hence $p_i\geq w_{ij}+p_j$. The length-zero path makes $p_i\geq0$. Conversely, iterating the inequalities along any path shows that every nonnegative feasible vector dominates its weight, and thus dominates $p$ coordinatewise. If all coordinates of $p$ were positive, subtracting their minimum would yield a smaller feasible nonnegative vector, contradicting leastness.
\end{proof}

\begin{theorem}[A quantitative identical-valuation neighborhood]
Let $n\geq2$. Suppose a normalized monotone function $v:2^M\to\mathbb R$ has marginal increments in $[0,1]$. Let $v_1,\ldots,v_n$ be any normalized valuations such that, for every bundle $S$,
\[
 |v_i(S)-v(S)|\leq\varepsilon\quad(i\in[n]).
\]
Then a complete allocation has envy-free nonnegative subsidies satisfying
\[
 \sum_i p_i\leq(n-1)\bigl(1+2(n-1)\varepsilon\bigr).
\]
In particular, identical normalized monotone valuations admit total subsidy at most $n-1$, and profiles with $\varepsilon\leq K/n$ admit total subsidy at most $(1+2K)(n-1)$.
\end{theorem}
\begin{proof}
Build $n$ bundles from empty bundles, adding each item to a bundle currently minimizing the common value $v$. Their common values always have spread at most one. Initially the spread is zero. If the current minimum is $a$, all values lie in $[a,a+1]$ by induction; monotonicity and the marginal bound keep the updated minimum bundle within that same interval. This proves the invariant.

Assign the resulting bundles to agents by a permutation maximizing the sum of actual values $\sum_i v_i(A_i)$. Such a permutation exists because there are finitely many. The associated weights have no positive cycle: cyclically assigning each agent the next bundle on a positive cycle would strictly improve that sum. Apply the lemma.

Write $e_i(S)=v_i(S)-v(S)$. On any simple path $i_0,\ldots,i_\ell$,
\begin{align*}
 \sum_{t=0}^{\ell-1}w_{i_t i_{t+1}}
 &=v(A_{i_\ell})-v(A_{i_0})
   +\sum_{t=0}^{\ell-1}\bigl(e_{i_t}(A_{i_{t+1}})-e_{i_t}(A_{i_t})\bigr)\\
 &\leq1+2\ell\varepsilon\leq1+2(n-1)\varepsilon.
\end{align*}
The common values telescope, which is the essential improvement over bounding each edge by one. The least subsidies therefore each satisfy this bound, and at least one is zero. Summing the other $n-1$ bounds proves the theorem.
\end{proof}

\begin{example}[The leading linear term cannot be removed]
One item with common value one forces total subsidy at least $n-1$: if $r$ receives it, envy-freeness requires $p_i\geq p_r+1$ for every $i\neq r$. The theorem's identical-valuation bound is therefore sharp.
\end{example}

\paragraph{Source relationship and remaining gap.}
Dupr\'e la Tour and Suzuki~\cite{source}, Section 6, explicitly ask whether a linear bound holds for general monotone valuations. Their full text was available for this note. The cycle-payment method is standard; the proof is included to expose the telescoping estimate. The neighborhood bound deteriorates quadratically when $\varepsilon$ is a fixed positive constant. Arbitrary heterogeneous valuations need not be close to any common function uniformly over all bundles, so the general linear target remains unresolved.

\begin{thebibliography}{9}
\bibitem{source} M. Dupr\'e la Tour and M. Suzuki.
\emph{Subquadratic Subsidies for Nonnegative or Nonpositive Valuations}, version 1, 2026, Sections 3 and 6.
\url{https://arxiv.org/html/2609.08272v1}.
\end{thebibliography}

\section{Completion Estimate}
15\%

\end{document}
